\chapter{How Evidence Forced the Idea of Atoms into Being}

\begin{kaichang}
Get a sugar cube from the kitchen (a small piece of rock sugar will do). Now do something apparently dull: break it in half, break one half in half again, and again, and again...

After only a few rounds, the pieces become too small to hold. Suppose you had infinite patience, infinitely skillful hands, and a microscope that could always magnify further. Could you continue forever? Can sugar be divided without end, always yielding ``smaller sugar''?

Before reading on, write down your answer: yes or no. People seriously debated this more than two thousand years ago. Its answer finally emerged through balances, measuring cylinders, and cold hard numbers, not a philosopher winning an argument. This chapter tells that story.
\end{kaichang}

\section{A Question Debated for Two Thousand Years}

The ancient Greek philosopher Democritus guessed that everything consisted of tiny, hard, indivisible particles. He called them atoms, from the Greek for ``uncuttable.'' In the same land, Aristotle's school held that matter was continuous and infinitely divisible, everything combining four qualities: earth, water, air, and fire.

Who was right? At the time, the debate could not be settled. Democritus's guess sounds closer to today's answer, but \textbf{being close to the answer is not the same as doing science}. He offered no evidence a skeptic could test and no prediction of what we should or should not observe if atoms existed. A conjecture can be beautiful yet impossible to confirm or refute. Testability is precisely the boundary between science and speculation. For the atomic idea to become science, someone had to corner it: explain the experimental numbers or leave.

That interrogation arrived around the turn of the nineteenth century. Its instrument was something you have at home: a scale.

\section{Three Unyielding Laws on the Balance}

Eighteenth-century chemists, notably Lavoisier, approached experiments as accountants: weigh everything before and after a reaction. Chemical change became an exercise in balancing books instead of watching a spectacle. Three laws emerged from the accounts.

\textbf{First: conservation of mass.} In a sealed vessel, a chemical reaction leaves total mass exactly unchanged. Burn a candle and its wax seems to disappear, but collect and weigh all the carbon dioxide and water vapor formed, and the books balance. Matter neither appears from nothing nor vanishes into nothing.

\textbf{Second: the law of definite proportions.} Whether water comes from a river or a laboratory, decomposition yields hydrogen and oxygen in the mass ratio $1:8$. You can check its stubbornness at home: boiling leaves white scale in the kettle, but both the remaining water and distilled water collected by condensing the vapor still decompose in the ratio $1:8$. Impurities may mingle, but water's recipe stays fixed. Copper oxide? However you prepare it, copper and oxygen have a mass ratio of about $4:1$. A pure compound's proportions are unchanging.

\textbf{Third: the law of multiple proportions.} Carbon and oxygen form two substances, carbon monoxide (the poisonous one) and carbon dioxide. For equal carbon masses, say \ud{12}{g}, the first contains \ud{16}{g} of oxygen and the second \ud{32}{g}---\textbf{exactly $1:2$}, a clean whole-number ratio. In several nitrogen--oxygen compounds, oxygen proportions follow $1:2:4$... Again, whole numbers.

Pause to appreciate how odd the third law is. If matter were continuous like dough and arbitrarily divisible, why would the same carbon insist on whole-number multiples of oxygen? It could just as well take $1.37$ or $2.61$ times as much. Yet nature favors integers. Imagine a rice shop selling only whole grains and never half a grain: you would immediately realize \textbf{rice comes in individual grains}.

\section{Dalton Brings Tiny Particles Back into Science}

Around 1803, the English teacher Dalton reasoned in just this way. Put the three laws together, he found, and one hypothesis makes them follow automatically:

\begin{itemize}[itemsep=2pt]
\item Each element consists of tiny particles, atoms, indivisible in chemical changes;
\item Atoms of the same element have identical masses and properties; those of different elements have different masses;
\item Compounds join different atoms in \textbf{simple whole-number ratios}. Reactions merely rearrange atoms, without adding, losing, creating, or destroying them.
\end{itemize}

Look back, and the three laws become consequences. If reactions neither create nor destroy atoms, mass must be conserved. Water molecules always combine one portion of hydrogen atoms with one portion of oxygen atoms, so composition is fixed. One carbon atom joins either one or two oxygen atoms, giving oxygen's $1:2$ ratio. Atoms combine whole, just as rice is sold by whole grains.

Notice the methodological turn. Democritus's atom was an ``I think.'' Dalton's was \textbf{a model that could balance the books}. It explained existing data and made predictions: if the law of multiple proportions failed for a compound, atomic theory would have to go. A theory becomes scientific when it openly states how it could be defeated.

Dalton's accounts were not all right at first, of course. Without the concept of molecular formulas, he thought water joined one hydrogen atom with one oxygen atom, equivalent to today's $\chem{HO}$, and wrote ammonia as $\chem{NH}$. What was missing? He knew mass ratios but not each atom's mass. Avogadro's hypothesis later closed the gap: equal gas volumes at equal temperature and pressure contain equal particle numbers. Comparing gas volumes before and after reaction reveals that water's hydrogen-to-oxygen ratio is $2:1$, not $1:1$. Even atomic theory's founder slipped on his own ladder. Models are corrected by successive data, not completed in one stroke.

\begin{qiaoliang}
Atomic theory arrived with a hair-raising implication: atom counts are absurdly large. Let us calculate, using the mole introduced in Chapter 5. A tiny drop of water weighs about \ud{0.05}{g}. How many molecules does it contain?
\[
\frac{\ud{0.05}{g}}{\ud{18}{g/mol}}\times \ud{6.02\times 10^{23}}{mol^{-1}} \approx 1.7\times 10^{21}\ \text{molecules}
\]
Seventeen quintillion---written out, 17 followed by 20 zeros. How large is that? Divide all the water in Earth's oceans, lakes, and rivers into cups of \ud{200}{mL}, and the count is on the order of $10^{21}$. In other words, \textbf{a tiny drop contains more water molecules than all Earth's water contains cupfuls}. Conversely, you can discover just how small each molecule is by doing the division yourself.
\end{qiaoliang}

\section{The Symbols Arrive Late}

Once atoms combine in integer ratios, bookkeeping symbols follow naturally. Dalton devised circle symbols: oxygen an empty circle, hydrogen a dotted circle, water two adjacent circles. It was as cumbersome as drawing comics. The Swedish chemist Berzelius later proposed letters from Latin names instead: O for one oxygen atom, H for one hydrogen atom, C for one carbon atom.

The real meaning of $\chem{H_2O}$ now becomes clear: \textbf{it is a census, not an abbreviation}. It reports that one water particle consists of 2 hydrogen atoms and 1 oxygen atom. The difference between $\chem{CO}$ and $\chem{CO_2}$ is not typography but the number of oxygen atoms accompanying carbon: $1:1$ versus $1:2$, the integer ratio in the law of multiple proportions. Symbols are particles' shadows. Particles come first, symbols second, which is why this book always leaves symbols until the end.

\section{Some Still Refused to Believe}

Dalton's theory explained everything, yet throughout the nineteenth century influential chemists and physicists, including the German chemist Ostwald, refused to accept it. Their objection was reasonable: atoms could not be seen or touched, so perhaps they were merely a bookkeeping trick that happened to work. Integer reaction ratios might be coincidences arising from some continuous law.

The debate lasted into the early twentieth century. A pollen particle's dance under a microscope ended it.

\begin{zhengju}
In 1827, the botanist Brown observed pollen particles suspended in water jiggling continuously and irregularly, as though invisible hands were pushing them. Lifeless powdered stone behaved the same, excluding living pollen as the explanation. For more than seventy years afterward, nobody could identify the source of the pushes.

In 1905, Einstein proposed a bold explanation: surrounding water molecules bombard pollen from all directions. The molecules are invisible, but their vast numbers of impacts cannot cancel perfectly at every instant. Sometimes one side hits harder, shoving the particle. He went further and calculated how far pollen should drift per unit time if water truly consisted of molecules of this size: a specific number measurable under a microscope.

The French physicist Perrin made that measurement. Over several years, he tracked suspended particles microscopically and found displacements matching Einstein's prediction, independently calculating particles per mole as well. The accounts balanced. After 1908, even Ostwald publicly conceded that atoms and molecules really existed. Remember the chain rather than a particular genius: an odd observation (jiggling pollen) $\rightarrow$ a testable quantitative explanation (molecular impacts) $\rightarrow$ a carefully designed measurement lasting years (Perrin) $\rightarrow$ the most stubborn skeptic changes his mind. Scientific consensus is forged this way, not by an authority's nod.
\end{zhengju}

\begin{shiyan}
\textbf{Observe ``Molecular Impacts'' at Home}

Materials: a glass of clean water, a drop of diluted red ink (or a little crushed watercolor-pen refill), white paper, and a magnifying glass (or a phone's macro mode).

Procedure: Gently place one ink drop on the surface of still water without stirring. Put the glass on white paper and watch for ten minutes.

What to observe: The ink does not spread evenly. Instead, thin, winding tendrils extend in irregular directions. These are traces left as randomly moving water molecules push dye particles.

Further exploration: Put one drop into hot water and another into cold water and compare simultaneously. Diffusion is visibly faster in hot water. Higher temperature means more vigorous molecular motion, a direct observation of temperature as molecular activity (Chapter 4 will return to this).

Safety: Use only clean water and stationery, with no open-flame heating. Clean up normally afterward.
\end{shiyan}

\section{The ``Uncuttable'' Particle Is Cut Open}

Dalton assumed atoms were indivisible, an assumption that never failed in chemical experiments. But at the nineteenth century's end, new toys---vacuum tubes, high voltage, and photographic plates---opened three windows into atoms.

\textbf{First window: the electron (1897).} Thomson studied cathode-ray tubes, evacuated glass tubes with high voltage across their ends. The cathode emits invisible rays that make glass fluoresce. Passing the rays through electric and magnetic fields, he found they always bent toward a positively charged plate, and their deflection was \textbf{entirely independent} of the gas inside or metal used for electrodes. The conclusion was firm: these were negative particles thousands of times lighter than hydrogen, the lightest atom, and \textbf{components shared by every atom}. Atoms contained smaller things. ``Uncuttable'' could no longer be taken literally.

\textbf{Second window: the nucleus (1911).} Thomson's picture was natural: an atom was a uniform positively charged pudding with electrons embedded like raisins. His student Rutherford tested it by bombarding very thin gold foil with fast positive particles (\chem{\alpha} particles, later identified as helium nuclei). If atoms were uniform puddings, these bulletlike particles should pass through almost undisturbed, and most did. But roughly one in twenty thousand \textbf{rebounded violently}, deflected through an astonishingly large angle. Rutherford later compared it to firing a shell at tissue paper and having it bounce back into your face.

Only one thing could reverse a fast particle: almost all atomic mass and all positive charge were crowded into a tiny core. An atom was not pudding but a mostly empty house, with a tiny heavy nucleus at its center and electrons active in the vast space outside. This conclusion convinced people partly because it gave quantitative, testable predictions: larger deflection angles should occur less often, by calculable amounts. Rutherford's assistant Geiger and student Marsden counted millions of flashes, and the numbers matched prediction point by point.

\textbf{Third window: protons and neutrons (1919 and 1932).} Rutherford later bombarded nitrogen with \chem{\alpha} particles and knocked out hydrogen nuclei---protons, basic nuclear constituents. But proton counts did not explain nuclear masses: a helium nucleus had 2 protons but nearly 4 protons' mass. Chadwick solved this in 1932 by discovering another nuclear particle, uncharged and almost as massive as a proton: the neutron. The family portrait was complete: protons and neutrons packed inside the nucleus, electrons outside. How proton number determines whether an atom is gold or iron is Chapter 7's story, the element's identity card.

\section{Three Models, Three Victories, Three Failures}

Pause to examine this relay of models: every stage was elegant, and every stage left a weakness.

\begin{itemize}[itemsep=2pt]
\item \textbf{Thomson's plum-pudding model}: explained electrical neutrality through balanced positive and negative charge, and why electrons could be pulled out. Its failure was uniformly distributed positive charge, directly demolished by the gold-foil experiment.
\item \textbf{Rutherford's nuclear model}: perfectly explained scattering and revealed the nucleus. Its problem was that contemporary electromagnetism predicted orbiting electrons would radiate energy and spiral into the nucleus almost instantly. Atoms should not survive a blink, yet they are stable. The model explained the nucleus's existence but not the atom's survival.
\item \textbf{Bohr's model (1913)}: Bohr stipulated that electrons occupy only certain allowed orbits, on which they \textbf{do not} radiate. Only jumps between orbits release the energy difference as light. This rescued atomic stability and explained hydrogen's discrete colored spectral lines (Chapter 9 develops the evidence). The failure was that these stipulations were patches: why those particular orbits? The model could not answer, and its calculations faltered for atoms more complex than hydrogen.
\end{itemize}

Do you see the pattern? Models were not simply wrong and discarded: each \textbf{explained one set of phenomena, then was cornered by the next}. Today's more accurate picture, electrons as probability clouds rather than little planets following fixed orbits, is Chapter 9's subject. For now, remember the rhythm: observations force a model, new observations force a new model. Nobody invented atoms in a single flash of inspiration; successive rounds of evidence forced the idea into being.

\section{How Small the Atom, How Empty Its Interior}

Put the scales side by side, and ``empty'' takes on a new meaning. Atomic diameters are on the order of $10^{-10}$ meters, varying somewhat by element; nuclear diameters are only about $10^{-15}$ to $10^{-14}$ meters. \textbf{A nucleus is only one hundred-thousandth of an atom's size.}

For a human-scale comparison, enlarge an atom to a standard soccer field. Its nucleus is a glass marble at the center, with only elusive electrons moving through the rest of the field. The apparently solid chair beneath you consists of matter that is more than $99.99999999999\%$ empty. You do not fall through because of electromagnetic repulsion between electrons, not because matter fills every space. Chapter 17 tells that story.

This also explains the gold-foil numbers. A nucleus occupies roughly $(10^{-14}/10^{-10})^2 = 10^{-8}$ of an atom's cross-sectional area. Even with over a thousand atomic layers, the probability of an \chem{\alpha} particle scoring a direct nuclear hit is only on the order of one in ten thousand, consistent with the observed one-in-twenty-thousand large-angle rebounds. A model earns trust when its numbers fit together.

\begin{wuqu}
``An atom is a solid little ball, like a miniature marble.'' This common mental picture is partly the fault of textbook illustrations. In reality, atomic mass crowds into a nucleus occupying less than a trillionth of the volume. Outside is no spherical shell, only probabilistically distributed electrons. Nor can any experiment inspect an atomic surface in the ordinary sense: its edge is inherently fuzzy. Imagining marbles prevents you from understanding chemical bonding (Chapter 17) or why most particles passed through gold foil. Solid balls in illustrations are a compromise imposed on the illustrator.
\end{wuqu}

\section{Where These Models Reach Their Limits}

As usual, let us state when this chapter's models work and when they fail:

\begin{itemize}[itemsep=2pt]
\item \textbf{``Atoms are indivisible''}: true in every chemical reaction---burning, rusting, and digestion merely rearrange atoms. It fails in nuclear reactions: nuclei undergo fission, fusion, and decay (Chapter 8), and protons and neutrons themselves contain more fundamental quarks. ``Uncuttable'' applies at chemistry's scale, not as an absolute truth.
\item \textbf{Bohr's model}: its strikingly accurate hydrogen spectrum makes it a useful introduction even today, but it fails quantitatively for many-electron atoms, and quantum mechanics has replaced the orbit concept itself. It is a ladder, not the destination.
\item \textbf{``Whole-number ratios''}: definite and multiple proportions hold for most compounds, but a few solids, including certain metal oxides, allow continuously varying composition within narrow limits. Atomic theory remains correct; these crystal recipes permit some missing atoms. Details belong to more advanced study.
\end{itemize}

\section{Look Again at the Sugar Cube}

Return to the opening question: can sugar be divided forever? You now have a full answer chain. Sugar contains carbon, hydrogen, and oxygen atoms in fixed proportions. Sucrose's $\chem{C_{12}H_{22}O_{11}}$ is a census: 12 carbon, 22 hydrogen, and 11 oxygen atoms per molecule. Repeated breaking makes grains smaller, but as long as it remains sugar, its smallest unit is one molecule of those 45 atoms. Divide further, and you no longer have smaller sugar but separated atoms: sugar disappears and its chemical properties are gone. Dalton's integer ratios and Perrin's particle counts guarantee this answer.

Everyday life reflects this chain. A rusty nail \textbf{gains} weight because oxygen atoms join from the air, not because more iron appears; weighing shows the books balance. A cut apple browns as molecules in the flesh and oxygen atoms rearrange. Burning sugar leaves black residue while its white sweetness disappears, a scene of molecules broken apart and atoms going separate ways. These visible changes all obey the same atomic ledger.

The half of Democritus's guess that was right---indivisible units exist---needs three corrections. The units are not the smallest things, because atoms contain protons, neutrons, and electrons. They are not hard solid balls, but spacious structures. And indivisibility holds only on the chemical-reaction stage.

\section{We Have Only Reached the Entrance Hall}

Atoms have long escaped textbooks to become real industries. The scanning tunneling microscope (STM), invented in 1981, scans a metal surface with a tip ending in a single atom. Tiny changes in current between tip and surface let it feel the arrangement of individual atoms. In 1990, IBM scientists moved xenon atoms one by one on nickel to spell their three-letter company logo. Chip engineers now design transistors in this atomic picture. Your phone fits over ten billion switches into a fingernail-sized chip because atoms are individual objects that can be arranged individually, the idea balances forced upon science two centuries ago.

This chapter has answered only whether atoms exist and roughly what lies inside. Beyond the entrance hall are more fascinating questions: why does the number of nuclear protons determine gold or iron, an inert gas or a volatile personality? The hundred-plus boxes on a chemistry classroom wall follow precisely this clue; Chapter 10 gives the whole map. Next, we obtain the identity card that determines an element: the number hidden in its nucleus.

\begin{tiaozhan}
We can roughly estimate Rutherford scattering with probability intuition. Imagine gold foil as a target wall of $N$ atomic layers. In each layer, the probability of an \chem{\alpha} particle hitting a nucleus is roughly the nuclear-to-atomic cross-sectional-area ratio: $(10^{-14}/10^{-10})^2 = 10^{-8}$. Rutherford's foil was a few hundred to over a thousand layers thick. Taking $N \approx 10^3$, the hit probability is about $10^3 \times 10^{-8} = 10^{-5}$, on the order of one in a hundred thousand. The observed large-angle fraction was about one in twenty thousand. Considering that a hit need not send a particle directly backward, these agree fairly well. We need not see atoms: counting projectiles and measuring rebounds reveals a target's internal structure, a fundamental skill of microscopic physics. Skipping this box does not interrupt the main story.
\end{tiaozhan}

\section*{Try It Yourself}

\begin{sikaoti}
\item Use Dalton's theory to explain why equal carbon amounts combining with \ud{16}{g} and \ud{32}{g} of oxygen provide strong evidence for atoms. What ratios would you expect if matter were continuous?
\item Thomson found the same cathode-ray deflection regardless of the tube's gas or electrode metal. Why is independence from material crucial evidence that electrons are components of all atoms?
\item Draw Rutherford scattering: one layer of gold atoms, each a large circle with a small nuclear dot, and five incoming \chem{\alpha} trajectories, four nearly straight and one sharply deflected. Note that the circles do not represent real atomic shells; they merely help us think.
\item If a hydrogen atom were enlarged to \ud{100}{m} across, about a playing field, how many millimeters across would its nucleus be? Explain why this scale ratio helps account for most particles passing straight through gold foil.
\item Bohr stipulated electrons occupy particular orbits and do not emit light while remaining there. What atomic property did this rescue? Why was the stipulation not a true explanation of nature?
\item Someone says, ``Ancient Greeks invented atoms, so science is just lucky guessing.'' Refute this in three sentences using the distinction between Democritus's conjecture and Dalton's atomic theory.
\end{sikaoti}
